% latexmk -pdf -pvc blatt3.tex
\documentclass{hott-übung}

\begin{document}

\setcounter{blattnummer}{3}
\setcounter{aufgabennummer}{0}

\blatt

\aufgabe{}
Sei $A:\mU$.
\begin{enumerate}[(a)]
\item Wenn $A$ eine Aussage ist und wir $a:A$ haben, dann ist $A$ kontrahierbar.
\item Wenn wir eine Funktion des Typs $A\to \isContr(A)$ haben, dann ist $A$ bereits eine Aussage. 
\end{enumerate}

\aufgabe{}
\begin{enumerate}[(a)]
\item Zeige $0\neq 1$, für die natürlichen Zahlen $0$ und $1$.  
\item Zeige $\neg(\eins\simeq \N)$.
\end{enumerate}

\aufgabe{}
\begin{enumerate}[(a)]
\item Seien $A$ ein Typ und $f:A\to A$ sodass $f\sim \id_A$. Zeige, dass $f$ eine Äquivalenz ist.
\item Seien $A,B,C$ Typen.
Konstruiere Abbildungen
\begin{align*}
\varphi : (A\times B \to C)\to (A\to B\to C)\quad\text{ und }  \quad\psi:(A\to B\to C)\to (A\times B\to C)
\end{align*}
zusammen mit $(f:A\times B \to C)\to \psi(\varphi(f))\sim f$ und $(g:A\to B\to C)\to \varphi(\psi(g))\sim g$.
\end{enumerate}

\aufgabe{2 aus 3}
Seien $f:A\to B$ und $g:B\to C$.
\begin{enumerate}[(a)]
\item Zeige, dass $g\circ f:A\to C$ eine Äquivalenz ist, wenn $f$ und $g$ Äquivalenzen sind.
\item Zeige, dass $f$ eine Äquivalenz ist, wenn $g$ und $g\circ f$ Äquivalenzen sind.
\item Zeige, dass $g$ eine Äquivalenz ist, wenn $f$ und $g\circ f$ Äquivalenzen sind.
\end{enumerate}

\bonus{}
Seien $A:\mU$ und $B:A\to\mU$. Zeige:
\begin{enumerate}[(a)]
\item Wenn $(x:A)\to \isContr(B(x))$, dann haben wir
  \[
    \left(\sum_{x:A}B(x)\right)\simeq A\rlap{.}
  \]
\item Wenn $A$ kontrahierbar ist, haben wir
  \[
    \left(\sum_{x:A}B(x)\right)\simeq B(\ast)\rlap{.}
  \]
\end{enumerate}

\end{document}
